\chapter{Requirements Analysis}\label{ch:requirements}

The canonical requirement register is version-controlled under \code{docs/requirements/}. Identifiers are
permanent: functional requirements FR-01 to FR-26, infrastructure requirements IR-01 to IR-09 and non-functional
requirements NFR-01 to NFR-16. This chapter summarises them and adds use-case identifiers (UC-01 to UC-13) for
the UML model in \autoref{ch:uml}. The verification status of each requirement, determined in the final
acceptance run of 9 October 2026, is given in \autoref{app:traceability}.

\section{Stakeholders and User Roles}

\begin{table}[htbp]
  \centering
  \caption{Stakeholders and roles}\label{tab:stakeholders}
  \begin{tabularx}{\textwidth}{@{}lL@{}}
    \toprule
    \textbf{Role} & \textbf{Interest and interaction} \\
    \midrule
    Researcher (user) & Imports and uploads papers, searches, asks questions, inspects citations, runs analyses. Single
      user; no accounts. \\
    Maintainer & Runs the evaluation harness and test suites, applies migrations, reads logs and decision records. \\
    Course evaluators & Assess requirements, design, verification evidence and documentation. \\
    arXiv (external system) & Supplies paper metadata and PDF files over HTTPS. \\
    Ollama (external system) & Serves the local language model used for answers and analyses. \\
    \bottomrule
  \end{tabularx}
\end{table}

\section{Functional Requirements}

\autoref{tab:fr} lists the functional requirements with the priority recorded in the register. ``Required''
requirements are mandatory for the MVP.

{\small
\begin{longtable}{@{}p{1.2cm}p{9.0cm}p{2.2cm}@{}}
  \caption{Functional requirements}\label{tab:fr} \\
  \toprule
  \textbf{ID} & \textbf{Requirement} & \textbf{Priority} \\
  \midrule
  \endfirsthead
  \toprule
  \textbf{ID} & \textbf{Requirement} & \textbf{Priority} \\
  \midrule
  \endhead
  FR-01 & Search the arXiv API by query or identifier and show metadata. & Required \\
  FR-02 & Import an arXiv paper; re-importing a known identifier updates rather than duplicates. & Required \\
  FR-03 & Upload a PDF with validation (type, size, pages, encryption, text layer) and safe storage. & Required \\
  FR-04 & Page-aware text extraction; every character traceable to a page and offset. & Required \\
  FR-05 & Deterministic chunking; chunks spanning pages record all pages. & Required \\
  FR-06 & Local embeddings with \code{all-MiniLM-L6-v2}; dimension validated as 384. & Required \\
  FR-07 & Persist papers, documents, chunks and embeddings in PostgreSQL with pgvector, constraints and migrations. & Required \\
  FR-08 & Semantic search returning ranked passages with score, chunk, paper and pages; corpus filter. & Required \\
  FR-09 & Grounded question answering with a local Ollama model, answering only from retrieved passages. & Required \\
  FR-10 & Citations resolved on the server; unresolvable citations flagged, never shown as valid. & Required \\
  FR-11 & Explicit insufficient-evidence response when context does not support an answer. & Required \\
  FR-12 & Single-paper summary with provenance. & Required \\
  FR-13 & Cross-paper synthesis with per-claim citations. & Required \\
  FR-14 & Structured extraction (task, method, dataset, metric, result, limitations), each field cited or \code{unknown}. & Required \\
  FR-15 & Multi-paper comparison table with caveats. & Required \\
  FR-16 & Corpus browsing: list, filter, details, delete with confirmation. & Required \\
  FR-17 & Ingestion job tracking with explicit states, actionable errors and retry. & Required \\
  FR-18 & Query history with answers, citations, configuration and latency. & Required \\
  FR-19 & Retrieval evaluation harness (Recall@$k$, MRR) on a versioned labelled set. & Required \\
  FR-20 & Answer-quality evaluation: citation validity, groundedness, abstention, latency. & Required \\
  FR-21 & Automated retrieval-quality gate in continuous integration. & Required \\
  FR-22 & Dashboard summarising corpus, jobs and recent questions. & Recommended \\
  FR-23 & Read-only settings view with models and dependency health. & Recommended \\
  FR-24 & Hybrid dense and lexical retrieval with rank fusion. & Recommended \\
  FR-25 & Cross-encoder re-ranking, adopted only if it improves the evaluation results. & Optional \\
  FR-26 & Optical character recognition for scanned PDFs. & Optional (deferred) \\
  \bottomrule
\end{longtable}
}

\section{Infrastructure Requirements}

IR-01 a versioned REST API (\code{/api/v1}) with typed schemas, a consistent error envelope and OpenAPI;
IR-02 environment-based configuration with no secrets in version control; IR-03 Docker Compose for the database,
backend and frontend; IR-04 continuous integration for lint, format, type checks, tests and build; IR-05 liveness
and per-dependency readiness endpoints; IR-06 structured JSON logs with request identifiers and no document text;
IR-07 an Ollama adapter with timeouts and error mapping; IR-08 background execution of ingestion jobs; IR-09
reproducible setup, troubleshooting and handover documentation.

\section{Non-Functional Requirements}

Non-functional requirements are targets to be measured; none is claimed without a recorded run.

\begin{table}[htbp]
  \centering
  \caption{Non-functional requirements}\label{tab:nfr}
  {\small
  \begin{tabularx}{\textwidth}{@{}llL@{}}
    \toprule
    \textbf{ID} & \textbf{Quality} & \textbf{Target} \\
    \midrule
    NFR-01 & Retrieval quality & Recall@5 $\geq 0.80$ on the pilot evaluation set \\
    NFR-02 & Retrieval quality & MRR reported for every retrieval configuration \\
    NFR-03 & Latency & 95th-percentile end-to-end question answering $\leq 3$\,s (hardware, model, warm/cold stated) \\
    NFR-04 & Latency & Search latency reported separately from generation latency \\
    NFR-05 & Availability & $\geq 99\%$ availability of search and Q\&A over a defined pilot window \\
    NFR-06 & Integrity & 100\% of displayed citations resolve to persisted chunks \\
    NFR-07 & Integrity & Abstention on unanswerable questions; precision and recall reported \\
    NFR-08 & Reproducibility & Same corpus and configuration give identical chunks and rankings \\
    NFR-09 & Reproducibility & Model, embedding model, chunk configuration, top-$k$ and prompt version recorded \\
    NFR-10 & Security & Upload limits, path-traversal protection, parameterised SQL, CORS, secrets, prompt injection \\
    NFR-11 & Privacy & No document text or secrets in logs; no data sent to hosted LLM services \\
    NFR-12 & Resource bounds & Upload size, pages, question length, top-$k$, tokens, context and concurrency bounded \\
    NFR-13 & Hardware fit & Runs on the reference laptop (16\,GB RAM, 4\,GB GPU, Windows 11) \\
    NFR-14 & Maintainability & CI green on every pull request \\
    NFR-15 & Accessibility & Keyboard navigation, labelled controls, usable at 360\,px width \\
    NFR-16 & Graceful degradation & Search works without Ollama; Q\&A returns an explicit 503 \\
    \bottomrule
  \end{tabularx}
  }
\end{table}

\section{Use Cases}

\autoref{tab:usecases} lists the use cases realised by the implementation. They are drawn in the use-case diagram
(\autoref{fig:uml-usecase}).

\begin{table}[htbp]
  \centering
  \caption{Use cases}\label{tab:usecases}
  {\small
  \begin{tabularx}{\textwidth}{@{}lLll@{}}
    \toprule
    \textbf{ID} & \textbf{Use case} & \textbf{Actors} & \textbf{Requirements} \\
    \midrule
    UC-01 & Discover papers on arXiv & Researcher, arXiv & FR-01 \\
    UC-02 & Import a paper from arXiv & Researcher, arXiv & FR-02, FR-04--FR-07 \\
    UC-03 & Upload a PDF & Researcher & FR-03--FR-07 \\
    UC-04 & Track an ingestion job (included by UC-02, UC-03) & Researcher & FR-17 \\
    UC-05 & Retry a failed job (extends UC-04) & Researcher, arXiv & FR-17 \\
    UC-06 & Browse and manage the corpus & Researcher & FR-16 \\
    UC-07 & Search passages & Researcher & FR-08, FR-24, FR-25 \\
    UC-08 & Ask a question & Researcher, Ollama & FR-09--FR-11 \\
    UC-09 & Inspect a citation (extends UC-08, UC-12) & Researcher & FR-10 \\
    UC-10 & Summarise or extract a paper & Researcher, Ollama & FR-12, FR-14 \\
    UC-11 & Synthesise and compare papers & Researcher, Ollama & FR-13, FR-15 \\
    UC-12 & View question history & Researcher & FR-18 \\
    UC-13 & View dashboard, settings and system status & Researcher & FR-22, FR-23, IR-05 \\
    \bottomrule
  \end{tabularx}
  }
\end{table}

The evaluation harness (FR-19 to FR-21) is a command-line tool used by the maintainer and in continuous integration;
it is not part of the user-facing system and is therefore not drawn as a use case.

\section{Assumptions and Constraints}

\begin{itemize}
  \item One user on one machine; all application ports bind to \code{127.0.0.1}; no authentication.
  \item No hosted language model, ever, as a fallback (contract rule; ADR-0001). If Ollama is unavailable, the
        system returns HTTP 503 with a reason.
  \item Python 3.12, managed by \code{uv} without changing the machine's Python 3.14, because PyTorch and some
        scientific packages often lag new Python releases on Windows (ADR-0004).
  \item Text-based PDFs only; scanned PDFs are rejected.
  \item The evaluation labels may be written by the AI assistant used by the team (ADR-0007); every item records its
        labeller and review status, and reports state how many items a person reviewed.
\end{itemize}

\section{Software and Hardware Requirements}

\begin{table}[htbp]
  \centering
  \caption{Reference environment}\label{tab:env}
  {\small
  \begin{tabularx}{\textwidth}{@{}lL@{}}
    \toprule
    \textbf{Item} & \textbf{Requirement / reference value} \\
    \midrule
    Hardware & AMD Ryzen 7 7435HS, 16\,GB RAM, NVIDIA RTX 3050 Laptop 4\,GB, Windows 11 \\
    Runtime & Docker Desktop (WSL2) with Compose; Ollama on the host with \code{qwen2.5:3b} \\
    Database & PostgreSQL 17 with pgvector 0.8.0 (container image \code{pgvector/pgvector:0.8.0-pg17}) \\
    Backend & Python 3.12, FastAPI, Pydantic v2, SQLAlchemy 2, Alembic, PyMuPDF, sentence-transformers (CPU PyTorch) \\
    Frontend & React, TypeScript, Vite, Vitest with Testing Library, oxlint; served by nginx 1.27 \\
    Tooling & Ruff, mypy (strict), pytest, GitHub Actions, pip-audit, npm audit \\
    \bottomrule
  \end{tabularx}
  }
\end{table}

\section{Acceptance Criteria}

The register defines numbered acceptance criteria (AC-$xx$.$n$) for each functional requirement in
\code{docs/requirements/acceptance-criteria.md}. The MVP was accepted on the following overall criterion: every
mandatory user journey (discovery, import, upload, provenance, embedding, search, grounded answering with verified
citations, abstention, analyses, corpus management, history, and graceful handling of invalid input and missing
services) works against the running application with real components, and no release-blocking defect remains. The
outcome (\autoref{ch:results}) was \emph{ready with documented limitations}.
